\begin{apartat}{1}
Un altaveu està format per un imant permanent en forma de disc i per una bobina per la qual circula un corrent elèctric. La bobina està unida a una membrana que participa dels moviments de la bobina.

\figura[0.45]{fig1.png}

\begin{itemize}[label=---]
  \item Com es mourà el conjunt bobina-membrana si fem circular un corrent continu per la bobina que, vist des de dalt, giri en sentit horari? \aladreta{\textit{[0,5 punts]}}
  \item Com es mourà el conjunt bobina-membrana si fem circular un corrent altern per la bobina? \aladreta{\textit{[0,5 punts]}}
\end{itemize}
Justifiqueu les respostes explicitant en cada cas la direcció i el sentit del camp magnètic produït per la bobina.
\end{apartat}

\begin{apartat}{1}
Necessitem més força sobre la bobina i per a aconseguir-ho cal que generi un camp magnètic més intens. Justifiqueu quin efecte tindria cada una de les modificacions següents sobre la intensitat del camp magnètic produït per la bobina:
\begin{itemize}[label=---]
  \item Un augment del nombre de voltes de la bobina. \aladreta{\textit{[0,5 punts]}}
  \item Un augment de la intensitat del corrent elèctric. \aladreta{\textit{[0,5 punts]}}
\end{itemize}
\end{apartat}

\nota{S'entén que en cada cas es manté constant el paràmetre que canvia en l'altra opció.}
